\chapter{Introduction}
% \label{chap:intro}


Attosecond optical spectroscopy has opened an unprecedented window into the real-time motion of electrons in matter, enabling direct observation of electronic excitations, charge migration, and correlation-driven dynamics on their natural timescale (10--1000~attoseconds). Over the past decade, techniques such as attosecond transient absorption spectroscopy (ATAS), high-harmonic spectroscopy (HHS), and time-resolved photoemission have revealed rich sub-femtosecond phenomena in atoms, molecules, and crystalline solids. However, despite this rapid progress, the attosecond response of \emph{disordered media} — including liquids, amorphous solids, polymers, and structurally complex hybrid materials — remains far less explored.

Disordered systems are ubiquitous in both natural and technological contexts. Their electronic structure is shaped not only by average bonding environments, but also by stochastic fluctuations of local potentials, short-range structural disorder
% enhanced scattering rates ?
and localization effects. These features strongly influence ultrafast electron dynamics: they modify mean free paths, induce rapid dephasing, and lead to spatially heterogeneous excitation pathways. Because many relevant photophysical and photochemical processes — photoinduced charge transfer, interfacial electron injection, initial steps of bond rearrangements — occur in such irregular environments, understanding attosecond-scale dynamics in disordered media is essential for accurate interpretation of measurements and rational materials design.

Recent experimental studies have demonstrated that attosecond techniques are indeed sensitive to disorder-induced scattering. High-harmonic spectroscopy in liquids has been shown to encode electron mean free paths and dephasing rates in the shape and phase of the emitted harmonics \cite{Mondal2023LiquidHHS}. Attosecond measurements in solids have captured ultrafast modifications of the band structure during strong-field excitation \cite{Schultze2014Si}. These results motivate the development of quantitative, predictive theoretical frameworks capable of describing attosecond processes in non-crystalline environments.

From a theoretical perspective, modeling attosecond optical responses of disordered systems requires a method that simultaneously accounts for (i) nonperturbative, real-time electronic dynamics and (ii) the explicit atomic-scale complexity of the medium. Real-time time-dependent density functional theory (RT-TDDFT) satisfies these requirements and has become a central tool for first-principles simulations of strong-field and attosecond phenomena in both molecules and solids. RT-TDDFT allows direct propagation of the electronic state in time, naturally capturing nonlinear responses, transient absorption signals, and high-harmonic emission. However, applying RT-TDDFT to disordered media remains challenging because of the need for large supercells, ensemble averaging over structural configurations, and the treatment of ultrafast scattering and decoherence — effects that are either absent or strongly suppressed in periodic crystals.

% Recent method-development papers \cite{SATO2021110274, Moitra2023CoreTDDFT} highlight the importance of improving exchange–correlation functionals, incorporating relativistic effects for core-level spectroscopy, and developing multiscale strategies that combine first-principles propagation with effective models for dissipation. However, a systematic and quantitative RT-TDDFT framework for attosecond spectroscopy of disordered materials has not been fully established.

\textbf{The objective of this thesis} is therefore to develop, validate, and apply a real-time TDDFT-based methodology for attosecond optical spectroscopy of disordered media. This includes constructing representative structural ensembles, implementing appropriate averaging procedures, computing ATAS and HHS observables, and identifying spectral signatures of disorder such as enhanced scattering,
% localization, 
ultrafast dephasing and relaxation. By comparing theoretical predictions with available experimental results, this work aims to clarify the microscopic origin of attosecond signals in non-crystalline environments and to provide practical guidelines for future attosecond spectroscopy of complex materials.





\section{Thesis Structure}
\label{sec:outline}

\textcolor{red}{TO DO}

